\section{Discussion}
\label{sec:discussion}

\subsection{Why $\alpha_c = 3.0$ is the right default}

The Pareto-frontier elbow at $\alpha_c = 3.0$ reflects a structural property
of the redistricting problem.  An intra-county edge in a well-drawn map
appears when the district boundary runs through a county interior --- a
configuration that the algorithm will avoid only if the cost of doing so
(measured in edge-cut increase) is exceeded by the stickiness penalty.

At $\alpha_c = 3.0$, the penalty is large enough to discourage unnecessary
intra-county cuts (those where an equally compact map could follow the
county line) but small enough to permit necessary cuts (those where
population balance requires crossing county lines).

The crossover between necessary and unnecessary cuts occurs at approximately
$\alpha_c = 3.0$ in our data, which can be understood geometrically: a
typical avoidable county split adds an edge cut equal to the county boundary
shared-length.  A stickiness multiplier of 3.0 makes this cut 3x more
expensive, which exceeds the 2.0--2.5x compactness benefit of the alternative
configuration in the vast majority of cases.

\subsection{Connection to legal requirements}

State constitutional provisions requiring ``preservation of political
subdivisions to the extent possible'' do not specify a quantitative
standard.  Courts have generally upheld maps with modest county-split
rates as satisfying this criterion, provided that splits are justified
by population-equality requirements~\citep{pildes2004democracy}.

The county-sticky approach provides a principled operationalisation of
``to the extent possible'': the algorithm preserves counties when doing so
does not compromise population balance or contiguity.
The DIA's specification of $\alpha_c = 3.0$ as the default implements
this principle algorithmically.

\paragraph{Constitutional criterion hierarchy.}
Some state constitutions require redistricting criteria to be applied
in strict priority order, with population equality paramount and
subdivision preservation secondary.
County-sticky edge weighting implements subdivision preservation as a
\emph{soft secondary criterion}: the population balance constraint
($\delta_{\max} \leq 0.5\%$) is enforced as a hard METIS parameter
(\texttt{ufactor = 5}), and county stickiness operates only within
the feasible region defined by that constraint.
This is consistent with the ``to the extent possible'' language in state
constitutional provisions, which is inherently bicriterion: subdivision
preservation is maximised subject to the overriding population-equality
requirement.

\paragraph{Partisan neutrality.}
County split reduction has no systematic partisan effect.
Comparing Democratic seat shares at $\alpha_c = 1.0$ (baseline) and
$\alpha_c = 3.0$ (county-sticky) across the 43 non-trivial states yields
a mean absolute seat difference of $|\Delta\text{seats}| = 0.3$.
This confirms that county preservation is geographically, not politically,
determined: the algorithm eliminates avoidable splits wherever they occur,
without regard to which party benefits from a given county boundary.

\paragraph{Transparency.}
The parameter $\alpha_c$ is published in the TIGER adjacency graph metadata,
making the weighting fully reproducible and auditable.  Any party can
independently verify that the county-sticky weights were applied correctly
by recomputing the adjacency graph with $\alpha_c = 3.0$.

\subsection{Limitations and open questions}

This study evaluates county-sticky weights only at census-tract resolution.
At block-group or block resolution, the geometry of county boundaries may
be captured more precisely, potentially changing the Pareto frontier.

The $\alpha_c$ sweep uses a fixed seed ($s_0$) per state.  Seed sensitivity
at $\alpha_c = 3.0$ has not been separately characterised; it is possible
that high-variance states (GA, NC from B.7) exhibit higher variance under
county-sticky weights.

Municipal boundaries (cities and towns) are not addressed in this paper.
Many state constitutions list both county and municipal preservation
as subdivision criteria.  Extending the county-sticky approach to
municipalities would require a two-level weight hierarchy:
$\alpha_c$ for county boundaries and $\alpha_m$ for municipal boundaries
within counties.  Preliminary experiments suggest $\alpha_m \approx 1.5$
is appropriate, but a systematic sweep remains to be done.
